\label{kap:zobrazeni}
\mathquote{Symmetry, as wide or as narrow as you may define its meaning, is one idea by which man through the ages has tried to comprehend and create order, beauty, and perfection.}{Hermann Weyl}

V předchozí kapitole jsme zkoumali útvary samotné. Teď se podíváme na to, jak s nimi hýbat: překlápět je, posouvat, otáčet a zvětšovat. Takovým operacím se říká geometrická zobrazení. Ve druhé části kapitoly pak využijeme všechno, co zatím o geometrii víme, k řešení konstrukčních úloh, tedy úloh, ve kterých máme útvar sestrojit jen pomocí pravítka a kružítka.

\section{Geometrická zobrazení}\label{sec:geom-zobrazeni}

\begin{definitionbox}
  \textbf{Zobrazení v rovině} je předpis, který každému bodu $X$ roviny přiřadí právě jeden bod $X'$ roviny. Bod $X$ se nazývá \textbf{vzor}, bod $X'$ \textbf{obraz}. Je to vlastně funkce z kapitoly~\ref{kap:funkce}, jejímž definičním oborem i oborem hodnot je rovina.

  \vspace{0.2cm}

  \textbf{Samodružný bod} je bod, který se zobrazí sám na sebe: $X' = X$. \textbf{Samodružný útvar} je útvar, který se zobrazí sám na sebe (jeho jednotlivé body ale samodružné být nemusí).
\end{definitionbox}

\section{Shodná zobrazení}\label{sec:shodna-zobrazeni}

\begin{definitionbox}
  \textbf{Shodné zobrazení} (shodnost) je zobrazení, které zachovává vzdálenosti: pro každé dva body $X, Y$ a jejich obrazy $X', Y'$ platí $|X'Y'| = |XY|$.

  \vspace{0.2cm}

  Shodné zobrazení převádí přímku na přímku, úsečku na stejně dlouhou úsečku, úhel na shodný úhel, rovnoběžky na rovnoběžky a každý útvar na útvar s ním shodný.
\end{definitionbox}

\subsection{Osová souměrnost}\label{sec:osova-soumernost}

\begin{definitionbox}
  \textbf{Osová souměrnost} $O(o)$ s osou $o$ přiřadí bodu $X\notin o$ bod $X'$ tak, že přímka $o$ je osou úsečky $XX'$. Body osy $o$ jsou samodružné.

  \vspace{0.2cm}

  Osová souměrnost \textbf{mění orientaci} útvarů: pokud jsou vrcholy trojúhelníku $ABC$ seřazeny proti směru hodinových ručiček, jsou vrcholy obrazu $A'B'C'$ po směru. Je to jako zrcadlo, ve kterém se levá ruka stane pravou.
\end{definitionbox}

\begin{center}
  \begin{tikzpicture}[scale=0.9]
    \draw[thick, colorofwisdom] (2,-0.5) -- (2,3.5) node[above] {\small $o$};
    \coordinate (A) at (0,0); \coordinate (B) at (1.4,0.6); \coordinate (C) at (0.4,2.2);
    \coordinate (A') at (4,0); \coordinate (B') at (2.6,0.6); \coordinate (C') at (3.6,2.2);
    \draw[thick] (A) -- (B) -- (C) -- cycle;
    \draw[thick] (A') -- (B') -- (C') -- cycle;
    \draw[dashed, gray] (A) -- (A') (B) -- (B') (C) -- (C');
    \node[below left] at (A) {\small $A$}; \node[below] at (B) {\small $B$}; \node[above] at (C) {\small $C$};
    \node[below right] at (A') {\small $A'$}; \node[below] at (B') {\small $B'$}; \node[above] at (C') {\small $C'$};
  \end{tikzpicture}
\end{center}

Útvar je \textbf{osově souměrný}, pokud existuje osová souměrnost, ve které je samodružný. Rovnostranný trojúhelník má $3$ osy souměrnosti, čtverec $4$, pravidelný $n$-úhelník $n$ a kružnice nekonečně mnoho (každá přímka procházející středem). Lidský obličej má přibližně jednu, ale jen přibližně.\footnote{Pokud chcete zažít opravdový horor, vyfoťte se zepředu a vytvořte si dvě fotky: jednu složenou ze dvou levých polovin obličeje, druhou ze dvou pravých. Ani jedna nebudete vy.}

\subsection{Středová souměrnost}

\begin{definitionbox}
  \textbf{Středová souměrnost} $S(S)$ se středem $S$ přiřadí bodu $X$ bod $X'$ tak, že $S$ je středem úsečky $XX'$. Jediným samodružným bodem je střed $S$. Samodružné jsou všechny přímky procházející bodem $S$.

  \vspace{0.2cm}

  Středová souměrnost \textbf{zachovává orientaci} a přímku zobrazí na rovnoběžnou přímku. Je to totéž co otočení o $180^\circ$.
\end{definitionbox}

\begin{center}
  \begin{tikzpicture}[scale=0.9]
    \coordinate (S) at (0,0);
    \coordinate (A) at (-2.4,-0.6); \coordinate (B) at (-1,-1.2); \coordinate (C) at (-1.5,0.5);
    \draw[thick] (A) -- (B) -- (C) -- cycle;
    \draw[thick, colorofwisdom] ($(S)!-1!(A)$) -- ($(S)!-1!(B)$) -- ($(S)!-1!(C)$) -- cycle;
    \draw[dashed, gray] (A) -- ($(S)!-1!(A)$) (B) -- ($(S)!-1!(B)$) (C) -- ($(S)!-1!(C)$);
    \filldraw (S) circle (1.5pt) node[below] {\small $S$};
    \node[left] at (A) {\small $A$}; \node[below] at (B) {\small $B$}; \node[above] at (C) {\small $C$};
    \node[right] at ($(S)!-1!(A)$) {\small $A'$}; \node[above] at ($(S)!-1!(B)$) {\small $B'$}; \node[below] at ($(S)!-1!(C)$) {\small $C'$};
    \node[font=\scriptsize, align=left, right] at (3,0) {$S$ je středem každé\\ úsečky $XX'$; trojúhelník\\ se `otočí vzhůru nohama'};
  \end{tikzpicture}
\end{center}

Středově souměrný je například rovnoběžník (střed je průsečík úhlopříček), kružnice, elipsa nebo graf liché funkce (střed v počátku, viz kapitola~\ref{kap:funkce}). Trojúhelník středově souměrný není nikdy.

\subsection{Posunutí a otočení}\label{sec:posunuti-otoceni}

\begin{definitionbox}
  \textbf{Posunutí} (translace) $T(\vec{u})$ o vektor $\vec{u}$ posune každý bod o tentýž vektor: $X' = X + \vec{u}$. Pokud $\vec{u}\neq\vec{0}$, nemá žádný samodružný bod. Samodružné jsou přímky rovnoběžné s vektorem $\vec{u}$.

  \vspace{0.2cm}

  \textbf{Otočení} (rotace) $R(S, \varphi)$ se středem $S$ o orientovaný úhel $\varphi$ přiřadí bodu $X\neq S$ bod $X'$ tak, že $|SX'| = |SX|$ a orientovaný úhel $XSX'$ má velikost $\varphi$. Jediný samodružný bod je střed $S$ (pokud $\varphi$ není násobek $360^\circ$).
\end{definitionbox}

\begin{center}
  \begin{tikzpicture}[scale=0.85]
    \begin{scope}
      \coordinate (A) at (0,0); \coordinate (B) at (1.2,0.2); \coordinate (C) at (0.5,1.2);
      \draw[thick] (A) -- (B) -- (C) -- cycle;
      \draw[thick] ($(A)+(2.5,0.8)$) -- ($(B)+(2.5,0.8)$) -- ($(C)+(2.5,0.8)$) -- cycle;
      \draw[->, thick, colorofwisdom] (0.5,-0.6) -- (3,0.2) node[midway, below right] {\small $\vec{u}$};
      \draw[dashed, gray] (C) -- ($(C)+(2.5,0.8)$);
      \node at (1.9,-1.4) {\small posunutí};
    \end{scope}
    \begin{scope}[xshift=7.5cm]
      \coordinate (S) at (0,0);
      \filldraw[colorofwisdom] (S) circle (2pt) node[below left] {\small $S$};
      \coordinate (A) at (1.5,0.2); \coordinate (B) at (2.5,0.4); \coordinate (C) at (1.9,1.2);
      \draw[thick] (A) -- (B) -- (C) -- cycle;
      \draw[thick, rotate around={70:(S)}] (A) -- (B) -- (C) -- cycle;
      \draw[dashed, gray] (S) -- (B);
      \draw[dashed, gray, rotate around={70:(S)}] (S) -- (B);
      \draw[->, colorofwisdom] (0.8,0.13) arc (9:79:0.8);
      \node[colorofwisdom] at (0.85,0.85) {\small $\varphi$};
      \node at (1.2,-1.4) {\small otočení};
    \end{scope}
  \end{tikzpicture}
\end{center}

\begin{note}
  Všechna shodná zobrazení v rovině jsou buď identita, osová souměrnost, posunutí, otočení, nebo tzv. posunuté zrcadlení (osová souměrnost složená s posunutím ve směru osy). Navíc každé z nich lze složit nejvýše ze tří osových souměrností. Osová souměrnost je tedy základní stavební kámen všech shodností:
  \begin{itemize}
    \item složením dvou osových souměrností s \textbf{rovnoběžnými} osami vzdálenými $d$ vznikne \textbf{posunutí} o $2d$ kolmo na osy,
    \item složením dvou osových souměrností s osami \textbf{různoběžnými} svírajícími úhel $\alpha$ vznikne \textbf{otočení} kolem jejich průsečíku o $2\alpha$.
  \end{itemize}
  Speciálně složení dvou souměrností s kolmými osami je otočení o $180^\circ$, tedy středová souměrnost. Proto když se podíváte do dvou zrcadel postavených do pravého úhlu, vidíte se tak, jak vás vidí ostatní, a ne zrcadlově.
\end{note}

\begin{center}
  \begin{tikzpicture}[scale=0.85]
    \draw[thick, colorofwisdom] (1,-0.6) -- (1,2.4) node[above] {\small $o_1$};
    \draw[thick, colorofwisdom] (2.5,-0.6) -- (2.5,2.4) node[above] {\small $o_2$};
    \coordinate (A) at (0,0); \coordinate (B) at (0.7,0.3); \coordinate (C) at (0.2,1.3);
    \draw[thick, fill=black!10] (A) -- (B) -- (C) -- cycle;
    \draw[thick, dashed, fill=black!5] (2,0) -- (1.3,0.3) -- (1.8,1.3) -- cycle;
    \draw[thick, fill=colorofwisdom!20] (3,0) -- (3.7,0.3) -- (3.2,1.3) -- cycle;
    \draw[->, gray] (0.2,1.6) -- (1.8,1.6); \draw[->, gray] (1.8,1.9) -- (3.2,1.9);
    \draw[<->] (1,-0.45) -- (2.5,-0.45) node[midway, below, font=\scriptsize] {$d$};
    \draw[->, very thick, colorofwisdom] (0,-1) -- (3,-1) node[midway, below, font=\scriptsize] {posunutí o $2d$};
  \end{tikzpicture}
  \hspace{1cm}
  \begin{tikzpicture}[scale=0.85]
    \coordinate (P) at (0,0);
    \draw[thick, colorofwisdom] (P) -- (0:3) node[right] {\small $o_1$};
    \draw[thick, colorofwisdom] (P) -- (35:3) node[above right] {\small $o_2$};
    \draw (0.8,0) arc (0:35:0.8); \node at (17:1.05) {\scriptsize $\alpha$};
    \coordinate (X) at (-20:2.2); \coordinate (X1) at (20:2.2); \coordinate (X2) at (50:2.2);
    \filldraw (X) circle (1.5pt) node[below right] {\small $X$};
    \filldraw[gray] (X1) circle (1.5pt) node[right] {\small $X_1$};
    \filldraw[colorofwisdom] (X2) circle (1.5pt) node[above] {\small $X'$};
    \draw[dashed, gray] (X) -- (X1) -- (X2);
    \draw[->, colorofwisdom] (-20:1.6) arc (-20:50:1.6);
    \node[font=\scriptsize, colorofwisdom] at (100:0.8) {otočení o $2\alpha$};
    \filldraw (P) circle (1.5pt);
  \end{tikzpicture}
\end{center}

\begin{problem}
  \begin{enumerate}[label=(\alph*)]
    \item Kolik os souměrnosti mají písmena A, H, N, O, S, X? Která z nich jsou středově souměrná?
    \item Určete samodružné body a samodružné přímky osové souměrnosti.
    \item Co vznikne složením dvou středových souměrností se středy $S_1\neq S_2$?
    \item Které pravidelné mnohoúhelníky jsou středově souměrné?
  \end{enumerate}
\end{problem}

\begin{solution}
  \begin{enumerate}[label=(\alph*)]
    \item A: $1$, H: $2$, N: $0$, O: $2$ (v typickém písmu, které není dokonale kulaté), S: $0$, X: $2$. Středově souměrná jsou H, N, O, S, X.
    \item Samodružné body: všechny body osy. Samodružné přímky: osa sama a všechny přímky na ni kolmé.
    \item Posunutí o vektor $2\overrightarrow{S_1S_2}$. Každá středová souměrnost je složením dvou osových s kolmými osami, a pokud osy zvolíme vhodně, dvě z nich se vyruší a zbydou dvě rovnoběžné osy.
    \item Ty se sudým počtem vrcholů.
  \end{enumerate}
\end{solution}

\section{Podobná zobrazení}\label{sec:podobna-zobrazeni}

\begin{definitionbox}
  \textbf{Podobné zobrazení} s koeficientem $k > 0$ je zobrazení, pro které platí $|X'Y'| = k\cdot|XY|$ pro každé dva body. Pro $k = 1$ je to shodnost, pro $k > 1$ zvětšení, pro $k < 1$ zmenšení. Podobné zobrazení zachovává úhly a poměry délek.

  \vspace{0.2cm}

  \textbf{Stejnolehlost} (homotetie) $H(S, \kappa)$ se středem $S$ a koeficientem $\kappa\neq 0$ přiřadí bodu $X$ bod $X'$ na přímce $SX$ takový, že
  \begin{equation*}
    \overrightarrow{SX'} = \kappa\cdot\overrightarrow{SX}.
  \end{equation*}
  \begin{itemize}
    \item Pro $\kappa > 0$ leží $X'$ na polopřímce $\overrightarrow{SX}$, pro $\kappa < 0$ na polopřímce opačné.
    \item Stejnolehlost je podobné zobrazení s koeficientem $k = |\kappa|$.
    \item Přímku zobrazí na \textbf{rovnoběžnou} přímku.
    \item Pro $\kappa = -1$ je to středová souměrnost.
  \end{itemize}
\end{definitionbox}

\begin{center}
  \begin{tikzpicture}[scale=0.8]
    \coordinate (S) at (0,0);
    \coordinate (A) at (2,0.3); \coordinate (B) at (3,1.2); \coordinate (C) at (1.6,1.6);
    \filldraw[colorofwisdom] (S) circle (2pt) node[left] {\small $S$};
    \draw[thick] (A) -- (B) -- (C) -- cycle;
    \draw[thick, colorofwisdom] ($(S)!2!(A)$) -- ($(S)!2!(B)$) -- ($(S)!2!(C)$) -- cycle;
    \draw[dashed, gray] (S) -- ($(S)!2!(A)$) (S) -- ($(S)!2!(B)$) (S) -- ($(S)!2!(C)$);
    \node[below] at (A) {\small $A$}; \node[right] at (B) {\small $B$}; \node[above left] at (C) {\small $C$};
    \node[below] at ($(S)!2!(A)$) {\small $A'$}; \node[right] at ($(S)!2!(B)$) {\small $B'$}; \node[above left] at ($(S)!2!(C)$) {\small $C'$};
    \node at (3,-1) {\small $H(S, 2)$};
  \end{tikzpicture}
\end{center}

\begin{example}
  Stejnolehlost se skvěle hodí ke konstrukci útvarů vepsaných do jiných útvarů. Úloha: do trojúhelníku $ABC$ vepište čtverec tak, aby jedna jeho strana ležela na straně $AB$ a zbylé dva vrcholy na stranách $AC$ a $BC$.

  Řešení: Nejdřív sestrojíme \textit{pomocný} čtverec $K_0L_0M_0N_0$, který má stranu $K_0L_0$ na $AB$ a vrchol $N_0$ na $AC$. Stačí zvolit libovolný bod $N_0$ na $AC$, spustit z něj kolmici na $AB$ a čtverec dokreslit. Jeho vrchol $M_0$ ale obecně na straně $BC$ ležet nebude.

  Teď si všimneme, že všechny takové čtverce (se stranou na $AB$ a vrcholem na $AC$) jsou navzájem stejnolehlé se středem $A$. Jejich vrcholy $M$ proto leží na jedné polopřímce $\overrightarrow{AM_0}$. Hledaný vrchol $M$ je tedy průsečík polopřímky $\overrightarrow{AM_0}$ se stranou $BC$. Z něj už čtverec dokreslíme snadno. Pomocný čtverec jsme `nafoukli' stejnolehlostí přesně na správnou velikost.
  \begin{center}
    \begin{tikzpicture}[scale=1.15]
      \coordinate (A) at (0,0); \coordinate (B) at (5,0); \coordinate (C) at (1.5,3);
      \draw[thick] (A) -- (B) -- (C) -- cycle;
      \draw[gray, dashed, fill=black!8] (0.5,0) rectangle (1.5,1);
      \node[font=\tiny, gray] at (1,0.5) {pomocný};
      \draw[colorofwisdom!70, ->] (A) -- (3.3,2.2);
      \fill[colorofwisdom!25] (0.9375,0) rectangle (2.8125,1.875);
      \draw[very thick, colorofwisdom] (0.9375,0) rectangle (2.8125,1.875);
      \filldraw (1.5,1) circle (1.2pt) node[right, font=\tiny] {$M_0$};
      \filldraw[colorofwisdom] (2.8125,1.875) circle (1.5pt) node[right, font=\scriptsize] {$M$};
      \node[below left] at (A) {\small $A$}; \node[below right] at (B) {\small $B$}; \node[above] at (C) {\small $C$};
    \end{tikzpicture}
  \end{center}
\end{example}

\section{Konstrukční úlohy}\label{sec:konstrukcni-ulohy}

Konstrukční úloha je úloha, ve které máme sestrojit útvar daných vlastností pouze pomocí pravítka (bez stupnice) a kružítka. Je to tradice, která sahá až k Eukleidovi. Proč jen pravítko a kružítko? Protože to jsou nástroje odpovídající dvěma nejjednodušším křivkám, přímce a kružnici. A protože Řekové měli rádi omezení.

\begin{definitionbox}
  \textbf{Řešení konstrukční úlohy} má čtyři části:
  \begin{enumerate}
    \item \textbf{Rozbor}: Předpokládáme, že úloha je vyřešená, nakreslíme si náčrtek a hledáme vztahy, ze kterých lze útvar sestrojit. Zvýrazníme zadané prvky.
    \item \textbf{Postup konstrukce}: Krok za krokem zapíšeme, jak útvar sestrojíme, ideálně symbolicky.
    \item \textbf{Konstrukce}: Samotné narýsování.
    \item \textbf{Diskuze}: Kolik má úloha řešení a za jakých podmínek? (Například trojúhelník ze stran $1$, $2$, $5$ neexistuje.)
  \end{enumerate}
\end{definitionbox}

\begin{definitionbox}
  \textbf{Symbolický zápis}: Používají se tyto zkratky:
  \begin{itemize}
    \item $AB$; $|AB| = 5\ \si{cm}$ \quad (sestrojení úsečky),
    \item $k$; $k(S; 3\ \si{cm})$ \quad (kružnice se středem $S$ a poloměrem $3\ \si{cm}$),
    \item $\overrightarrow{AX}$; $|\angle BAX| = 60^\circ$ \quad (polopřímka svírající s $AB$ úhel $60^\circ$),
    \item $C$; $C\in k\cap\overrightarrow{AX}$ \quad (bod jako průsečík),
    \item $\triangle ABC$ \quad (výsledný trojúhelník).
  \end{itemize}
\end{definitionbox}

\subsection{Množiny bodů dané vlastnosti}\label{sec:mnoziny-bodu}

Jádrem většiny konstrukcí je hledání bodu, který má splňovat dvě podmínky zároveň. Každá podmínka určí nějakou množinu bodů a hledaný bod leží v jejich průniku. Opět množiny, opět průnik, opět kapitola 2.

\begin{definitionbox}
  \begin{center}
    \renewcommand{\arraystretch}{1.4}
    \begin{tabular}{|p{6.5cm}|p{6.5cm}|}
      \hline
      Množina bodů, které\ldots & je \\
      \hline
      mají od bodu $S$ vzdálenost $r$ & kružnice $k(S; r)$ \\
      mají stejnou vzdálenost od bodů $A$ a $B$ & osa úsečky $AB$ \\
      mají od přímky $p$ vzdálenost $d$ & dvojice rovnoběžek s $p$ ve vzdálenosti $d$ \\
      mají stejnou vzdálenost od dvou různoběžek & dvojice os úhlů, které přímky svírají \\
      z nichž je úsečka $AB$ vidět pod pravým úhlem & Thaletova kružnice nad $AB$ (bez bodů $A$, $B$) \\
      z nichž je úsečka $AB$ vidět pod úhlem $\alpha$ & dva kruhové oblouky nad $AB$ souměrné podle $AB$ (podle věty o obvodovém úhlu) \\
      \hline
    \end{tabular}
  \end{center}
\end{definitionbox}

\begin{center}
  \begin{tikzpicture}[scale=0.62]
    \begin{scope}
      \coordinate (A) at (0,0); \coordinate (B) at (2,0);
      \draw[thick] (A) -- (B); \filldraw (A) circle (1.5pt) node[below] {\tiny $A$}; \filldraw (B) circle (1.5pt) node[below] {\tiny $B$};
      \draw[thick, colorofwisdom] (1,-1.3) -- (1,1.7);
      \filldraw[colorofwisdom] (1,1.1) circle (1.5pt); \draw[dashed, gray] (A) -- (1,1.1) -- (B);
      \node[font=\scriptsize, align=center] at (1,-1.9) {osa úsečky};
    \end{scope}
    \begin{scope}[xshift=4cm]
      \draw[thick] (0,0) -- (2.4,0) (0,0) -- (50:2.4);
      \draw[thick, colorofwisdom] (0,0) -- (25:2.5);
      \coordinate (P) at (25:1.6); \filldraw[colorofwisdom] (P) circle (1.5pt);
      \draw[dashed, gray] (P) -- ({1.6*cos(25)},0) (P) -- ($(0,0)!(P)!(50:2.4)$);
      \node[font=\scriptsize] at (1.1,-1.9) {osa úhlu};
    \end{scope}
    \begin{scope}[xshift=8cm]
      \draw[thick] (0,0) -- (2.6,0) node[right, font=\tiny] {$p$};
      \draw[thick, colorofwisdom] (0,0.9) -- (2.6,0.9) (0,-0.9) -- (2.6,-0.9);
      \draw[<->, gray] (1.3,0) -- (1.3,0.9) node[midway, right, font=\tiny] {$d$};
      \node[font=\scriptsize] at (1.3,-1.9) {rovnoběžky ve vzd. $d$};
    \end{scope}
    \begin{scope}[xshift=13cm]
      \coordinate (A) at (-1.2,0); \coordinate (B) at (1.2,0); \coordinate (V) at (115:1.2);
      \draw[thick, colorofwisdom] (0,0) circle (1.2);
      \draw[thick] (A) -- (B); \draw (A) -- (V) -- (B);
      \filldraw (A) circle (1.5pt) (B) circle (1.5pt);
      \node[font=\scriptsize] at (0,-1.9) {Thaletova kružnice};
    \end{scope}
  \end{tikzpicture}
\end{center}

\begin{example}
  Sestrojme trojúhelník $ABC$, je-li dáno $c = 6\ \si{cm}$, $v_c = 3\ \si{cm}$ a $\gamma = 90^\circ$.

  \textbf{Rozbor}: Úsečku $AB$ sestrojíme rovnou. Bod $C$ musí splňovat dvě podmínky. Úhel $ACB$ je pravý, takže $C$ leží na Thaletově kružnici nad $AB$. Vzdálenost $C$ od přímky $AB$ je $v_c$, takže $C$ leží na rovnoběžce s $AB$ ve vzdálenosti $3\ \si{cm}$.

  \textbf{Postup}:
  \begin{enumerate}
    \item $AB$; $|AB| = 6\ \si{cm}$,
    \item $S$; $S$ je střed $AB$,
    \item $k$; $k(S; 3\ \si{cm})$ (Thaletova kružnice),
    \item $p$; $p\parallel AB$, $|pAB| = 3\ \si{cm}$,
    \item $C$; $C\in k\cap p$,
    \item $\triangle ABC$.
  \end{enumerate}

  \textbf{Diskuze}: Thaletova kružnice má poloměr $3\ \si{cm}$, tedy přesně $v_c$. Rovnoběžka $p$ je její tečnou a má s ní jediný společný bod. V jedné polorovině je tedy jedno řešení (v druhé polorovině je shodné, souměrné podle $AB$). Pokud by bylo $v_c > \frac{c}{2}$, úloha by neměla řešení. Pokud $v_c < \frac{c}{2}$, má v každé polorovině dvě (navzájem souměrná) řešení.
  \begin{center}
    \begin{tikzpicture}[scale=0.6]
      \coordinate (A) at (0,0); \coordinate (B) at (6,0); \coordinate (S) at (3,0); \coordinate (C) at (3,3);
      \draw[dashed, colorofwisdom] (S) circle (3);
      \draw[colorofwisdom!70] (-1,3) -- (7,3) node[right, font=\scriptsize] {$p$};
      \fill[colorofwisdom!15] (A) -- (B) -- (C) -- cycle;
      \draw[thick] (A) -- (B) -- (C) -- cycle;
      \draw ($(C)!0.12!(A)$) -- ($(C)!0.12!(A)!0.12!(B)$);
      \filldraw (S) circle (1.5pt) node[below] {\scriptsize $S$};
      \node[below left] at (A) {\small $A$}; \node[below right] at (B) {\small $B$}; \node[above] at (C) {\small $C$};
      \draw[<->, gray] (6.5,0) -- (6.5,3) node[midway, right, font=\scriptsize] {$v_c$};
      \node[font=\scriptsize, colorofwisdom] at (0.4,2.7) {$k$};
    \end{tikzpicture}
  \end{center}
\end{example}

\begin{example}\label{ex:tecny-z-bodu}
  Sestrojme tečny z bodu $P$ ke kružnici $k(S; r)$, kde $P$ leží vně kružnice.

  \textbf{Rozbor}: Tečna je kolmá na poloměr v bodě dotyku $T$. Úhel $STP$ je tedy pravý a bod $T$ leží na Thaletově kružnici nad $SP$.

  \textbf{Postup}:
  \begin{enumerate}
    \item $O$; $O$ je střed $SP$,
    \item $l$; $l(O; |OS|)$,
    \item $T_1, T_2$; $T_1, T_2\in k\cap l$,
    \item $t_1, t_2$; $t_1 = PT_1$, $t_2 = PT_2$.
  \end{enumerate}
  Úloha má vždy dvě řešení, pokud $P$ leží vně kružnice.
  \begin{center}
    \begin{tikzpicture}[scale=0.9]
      \coordinate (S) at (0,0); \coordinate (P) at (5,0); \coordinate (O) at (2.5,0);
      \draw[thick] (S) circle (1.5) node[above left, xshift=-0.9cm, yshift=0.9cm] {\small $k$};
      \draw[dashed, colorofwisdom] (O) circle (2.5);
      \filldraw (S) circle (1.5pt) node[below] {\small $S$};
      \filldraw (P) circle (1.5pt) node[right] {\small $P$};
      \filldraw (O) circle (1.5pt) node[below] {\small $O$};
      \coordinate (T1) at (0.45,1.431); \coordinate (T2) at (0.45,-1.431);
      \filldraw (T1) circle (1.5pt) node[above left] {\small $T_1$};
      \filldraw (T2) circle (1.5pt) node[below left] {\small $T_2$};
      \draw[thick] ($(P)!1.25!(T1)$) -- (P) -- ($(P)!1.25!(T2)$);
      \draw[gray] (S) -- (T1) (S) -- (T2) (S) -- (P);
      \node[colorofwisdom] at (2.5,2.75) {\small $l$};
    \end{tikzpicture}
  \end{center}
\end{example}

\begin{example}
  Sestrojme trojúhelník $ABC$, je-li dáno $c = 5\ \si{cm}$, $\gamma = 60^\circ$ a $t_c = 4\ \si{cm}$ (těžnice z vrcholu $C$).

  \textbf{Rozbor}: Úsečku $AB$ sestrojíme. Těžnice vede ze středu $S_c$ strany $AB$ do bodu $C$, takže $C$ leží na kružnici $k(S_c; 4\ \si{cm})$. Z bodu $C$ je úsečka $AB$ vidět pod úhlem $60^\circ$, takže $C$ leží na jednom ze dvou kruhových oblouků, ze kterých je $AB$ vidět pod úhlem $60^\circ$. Střed $O$ takového oblouku leží na ose $AB$ a středový úhel $AOB$ je dvojnásobkem obvodového, tedy $120^\circ$.

  \textbf{Postup} (v jedné polorovině):
  \begin{enumerate}
    \item $AB$; $|AB| = 5\ \si{cm}$,
    \item $S_c$; $S_c$ je střed $AB$,
    \item $o$; $o$ je osa $AB$,
    \item $O$; $O\in o$, $|\angle AOB| = 120^\circ$ (stačí sestrojit $|\angle OAB| = 30^\circ$),
    \item $m$; $m(O; |OA|)$,
    \item $k$; $k(S_c; 4\ \si{cm})$,
    \item $C$; $C\in m\cap k$,
    \item $\triangle ABC$.
  \end{enumerate}
  \textbf{Diskuze}: Kružnice $k$ a oblouk $m$ se mohou protnout ve dvou, jednom nebo žádném bodě. V našem případě dostaneme v jedné polorovině dva průsečíky, které jsou souměrné podle osy $o$, takže vzniknou dva shodné trojúhelníky. Až na shodnost je tedy řešení jedno.
  \begin{center}
    \begin{tikzpicture}[scale=0.75]
      \coordinate (A) at (0,0); \coordinate (B) at (5,0); \coordinate (Sc) at (2.5,0); \coordinate (O) at (2.5,1.443);
      \draw[thick, colorofwisdom] (A) arc (210:-30:2.887);
      \draw[dashed] (Sc) circle (4);
      \draw[gray] (2.5,-1) -- (2.5,4.8) node[above, font=\scriptsize] {$o$};
      \draw[gray] (A) -- (O) -- (B);
      \draw (O) ++(210:0.45) arc (210:330:0.45); \node[font=\tiny] at (2.5,0.75) {$120^\circ$};
      \coordinate (C1) at (4.642,3.378); \coordinate (C2) at (0.358,3.378);
      \fill[colorofwisdom!15] (A) -- (B) -- (C1) -- cycle;
      \draw[thick] (A) -- (B) -- (C1) -- cycle;
      \draw[gray] (Sc) -- (C1) node[midway, right, font=\scriptsize] {$t_c$};
      \filldraw (O) circle (1.5pt) node[left, font=\scriptsize] {$O$};
      \filldraw (Sc) circle (1.5pt) node[below, font=\scriptsize] {$S_c$};
      \filldraw (C1) circle (1.5pt) node[right, font=\scriptsize] {$C$};
      \filldraw[gray] (C2) circle (1.5pt) node[left, font=\scriptsize] {$C'$};
      \node[below left] at (A) {\small $A$}; \node[below right] at (B) {\small $B$};
      \node[colorofwisdom, font=\scriptsize] at (4.6,2.1) {$m$}; \node[font=\scriptsize] at (6.6,1) {$k$};
    \end{tikzpicture}
  \end{center}
\end{example}

\begin{problem}
  Proveďte rozbor a napište postup konstrukce:
  \begin{enumerate}[label=(\alph*)]
    \item trojúhelník $ABC$: $a = 5\ \si{cm}$, $b = 4\ \si{cm}$, $\gamma = 45^\circ$,
    \item trojúhelník $ABC$: $c = 7\ \si{cm}$, $v_c = 3\ \si{cm}$, $t_c = 4\ \si{cm}$,
    \item kružnice o poloměru $2\ \si{cm}$, která se dotýká dvou daných různoběžek,
    \item kosočtverec, je-li dána jeho strana $a = 4\ \si{cm}$ a úhlopříčka $u = 6\ \si{cm}$.
  \end{enumerate}
\end{problem}

\begin{solution}
  \begin{enumerate}[label=(\alph*)]
    \item Věta sus: $CB$; $|CB| = 5$; $\overrightarrow{CX}$; $|\angle BCX| = 45^\circ$; $A$; $A\in\overrightarrow{CX}$, $|CA| = 4$; $\triangle ABC$. Jedno řešení.
    \item $AB$; $|AB| = 7$; $S_c$ střed $AB$; $p\parallel AB$ ve vzdálenosti $3$; $k(S_c; 4)$; $C\in p\cap k$. Protože $t_c > v_c$, kružnice protne rovnoběžku ve dvou bodech, které dají dva shodné trojúhelníky (souměrné podle osy $AB$).
    \item Střed hledané kružnice má od obou přímek vzdálenost $2$, leží tedy na průsečíku rovnoběžek s danými přímkami ve vzdálenosti $2$ (což je zároveň na osách úhlů). Takové body jsou čtyři, úloha má čtyři řešení.
    \item Úhlopříčky kosočtverce se navzájem kolmo půlí. Sestrojíme úhlopříčku $AC$ délky $6$, její osu a na ní body $B, D$ ve vzdálenosti $4$ od $A$ (průsečíky kružnice $k(A; 4)$ s osou). Řešení existuje, protože $4 > 3$.
  \end{enumerate}
\end{solution}

\section{Odbočka --- Tři nemožné úlohy}\label{sec:nemozne-ulohy}

Řekové nám zanechali tři konstrukční úlohy, které nedokázali vyřešit, a na kterých si pak dva tisíce let lámaly zuby generace matematiků i amatérů:
\begin{enumerate}
  \item \textbf{Trisekce úhlu}: rozdělit libovolný úhel na tři shodné části.
  \item \textbf{Zdvojení krychle}: sestrojit hranu krychle, která má dvojnásobný objem než daná krychle. Neboli sestrojit délku $\sqrt[3]{2}$.
  \item \textbf{Kvadratura kruhu}: sestrojit čtverec se stejným obsahem jako daný kruh. Neboli sestrojit délku $\sqrt{\pi}$.
\end{enumerate}

Ke zdvojení krychle se váže temná legenda. Kolem roku 430 před naším letopočtem zachvátil Athény mor, na který zemřela asi čtvrtina obyvatel, včetně státníka Perikla. Athéňané poslali posly k věštírně na ostrov Délos, aby se zeptali, jak mor zastavit. Odpověď zněla: zdvojnásobte Apollónův oltář. Oltář měl tvar krychle. Athéňané nejdřív zdvojnásobili jeho hranu, ale tím objem vzrostl osmkrát a mor neustal. Úloze se proto dodnes říká \textbf{délský problém}. Podle Plútarcha pak Platón vysvětlil, že bůh oltář ve skutečnosti nepotřebuje. Chtěl jen, aby se Řekové víc věnovali geometrii.

Odpověď přišla až v 19. století a zněla: \textit{nejde to}. Pomocí pravítka a kružítka lze sestrojit jen délky, které vzniknou z daných délek sčítáním, odčítáním, násobením, dělením a druhou odmocninou. Číslo $\sqrt[3]{2}$ takhle získat nejde, a stejně tak (u obecného úhlu) trisekce vede na kubickou rovnici, kterou druhými odmocninami vyřešit nejde. Obojí dokázal v roce 1837 francouzský matematik Pierre Wantzel. Kvadraturu kruhu pohřbil v roce 1882 Ferdinand von Lindemann důkazem, že $\pi$ je transcendentní, tedy není kořenem žádného mnohočlenu s celočíselnými koeficienty.

Wantzel sám zemřel ve 33 letech. Jeho kolega Saint-Venant o něm napsal, že pracoval po večerech, málo spal, střídavě se uchyloval ke kávě a opiu a jedl, kdy si vzpomněl. A přestože je kvadratura kruhu dokázaně nemožná už přes sto let, matematické ústavy po celém světě dodnes dostávají dopisy od lidí, kteří ji `vyřešili'. Výraz `kvadratura kruhu' se v mnoha jazycích dokonce ujal jako synonym pro nemožný úkol.

\section{Shrnutí a otázky k ústní zkoušce}

\begin{problem}
  Odpovězte vlastními slovy:
  \begin{enumerate}
    \item Co je geometrické zobrazení? Co je samodružný bod a samodružný útvar?
    \item Co je shodné zobrazení? Popište osovou a středovou souměrnost, posunutí a otočení. Jaké mají samodružné body?
    \item Která shodná zobrazení mění orientaci?
    \item Co vznikne složením dvou osových souměrností?
    \item Co je podobné zobrazení a stejnolehlost? Jaké vlastnosti má stejnolehlost?
    \item Uveďte příklady osově a středově souměrných útvarů.
    \item Z jakých částí se skládá řešení konstrukční úlohy?
    \item Jaké znáte množiny bodů dané vlastnosti a jak se využívají v konstrukčních úlohách?
    \item Jak sestrojíte tečnu ke kružnici z bodu ležícího vně kružnice?
    \item Které klasické konstrukční úlohy nelze vyřešit pravítkem a kružítkem a proč?
  \end{enumerate}
\end{problem}
